\section{Score-P} 
\label{section:localisation:scorep}

We here give an overview of getting started with building profiles of parallel
applications using Score-P which is a performance infrastructure. We here will
look briefly at using the Score-P compiler wrappers to build instrumented
binaries which, when executed, produce profiles that can be visualised with the
Cube GUI. It should be noted that Score-P offers a user instrumentation API to
target specific kernels and functions, but here we focus on using the Score-P
compiler wrappers to instrument the entire binary.

An individual Score-P install is built with a specific combination of compiler
and MPI library, meaning on clusters multiple Score-P installs are often
required. It is best to reach out to cluster support teams for help with
getting the relevant combinations of compiler and Score-P as it might be useful
to compare compilers during analysis.

It should be noted that as with many profiling and performance tools, an
application's runtime can increase (potentially quite significantly). Hence, it
is always advised to run the application without any tooling to note the
runtime and also resource usage such as maximum memory footprint.

\subsection{Instrumentation}

We load our installation of Score-P and then check that we have the Score-P
compiler wrappers typically of the form \texttt{scorep-x} for a given compiler,
\texttt{x}. If building with a build system like CMake it is best to set the
\texttt{SCOREP\_WRAPPER=off} environment variable to stop instrumentation during
configuration.

This compilation will then have produced an instrumented binary which we can
run in a job script as we before (making sure to load in the Score-P library),
though there is some good practice to consider first. As with tools like MAQAO,
Score-P will write output files to an experiment directory which will by
default be given a long name based on the time and date. Instead, we can be
more exact and set an experiment directory name which is more meaningful with
the \texttt{export SCOREP\_EXPERIMENT\_DIRECTORY=...} environment variable.

Once the code is finished, there will be a local experimental directory which
will contain several files such as the files ending with \texttt{*.cubex} which
can be visualised in the Cube GUI tool. However, we can also do some further
post-processing with the \texttt{square -s <experiment-directory>}, and then
visualise with the CubeGUI. For more information in using the CubeGUI please
refer to the documentation in Ref~\cite{}. Here we are mainly focused on a
workflow of generating data in Score-P, we discuss data interpretation
elsewhere in the document. Next, we touch on several configurations to gather
more low-level information during the profiling.

\subsubsection{Performance counters}

Many performance tools allow the user to gather information on hardware
performance counters. With Score-P we can, for example, set the
\begin{verbatim}
export SCOREP_METRIC_PAPI=PAPI_TOT_INS,PAPI_FP_OPS
\end{verbatim}
environment variable which will gather hardware counters via PAPI, where here
we have selected the total instructions and floating point operations. The
availability of PAPI counters can be found by the \texttt{papi\_avail} command. 

We can also use performance metrics from Linux \texttt{perf}, but we refer the
reader to the Score-P documentation for further information.

These performance counters then appear in the CubeGUI as further metrics which
we can breakdown across the further dimensions. However, as always it can be
easy to generate a lot of data; performance counters can be very powerful but
it useful to have an aim in collecting them, to guide an analysis.

\subsubsection{Memory}

Beyond hardware performance counters, Score-P also supports memory based
metrics through the
\begin{verbatim}
export SCOREP_MEMORY_RECORDING=true
export SCOREP_MPI_MEMORY_RECORDING=true
\end{verbatim}
environment variables which can help detecting memory related issues such as
memory leaks. The support for these features can vary across languages and, as
with performance counters, it is useful to keep in mind that gathering this
data should only be carried out if there is an aim for the analysis.

\subsection{Filtering}

A crucial aspect of using tools like Score-P is filtering. Running instrumented
binaries can show large increases in runtime, and in the next section we can
see how to use the Score-P infrastructure to produce trace analyses which can
generate very large amounts of data. Hence, filtering is a major step in this
workflow.

To begin with we can use the \texttt{scorep-score} command to generate a
summary of the experiment directory as below
\begin{verbatim}
scorep-score ./profile.cubex
\end{verbatim}
Running this command on the Score-P generated profile we get predictions of how
large a trace analysis would be with the current setup and often these can be
tens of GigaBytes or even TeraBytes. Crucially we will also see information
regarding the total memory required by Score-P, which can be set by the
\texttt{SCOREP\_TOTAL\_MEMORY} environment vaariable.

We can also recover a more detailed
breakdown by passing the \texttt{-r} flag to \texttt{scorep-score} which then
shows lots of information including the number of times each function is
called. This data is very valuable in building our filtering but first we need
an initial filter file.

We can generate an initial filter file by using the command
\begin{verbatim}
score-p -g ./profile.cubex
\end{verbatim}
Then we can factor this into our prediction for the trace analysis size with
the \texttt{-f} flag
\begin{verbatim}
scorep-score -f initial-scorep.filter -r ./profile.cubex
\end{verbatim}
This should show a reduction in the predicted trace size, although we can
typically do much better by using the outputs of the \texttt{scorep-score -r}
command to iteratively add further functions and files to our filtering file.

Performance counter data also adds to the size of any resulting trace analyses,
and their effects can be included in the size predictions with the \texttt{-c}
flag followed by the number of performance counters. For example, in this code
snippet
\begin{verbatim}
iscorep-score -f scorep.filter -c 3 -r ./profile.cubex
\end{verbatim}
we are predicting to gather 3 performance counters.

Once we decided on a filtering file, it is included with
\texttt{SCOREP\_FILTERING\_FILE} environment variable. Now that we have built a
filtering system with Score-P, in the next section we will demonstrate how the
Score-P infrastructure can be used to build automated trace analyses with
Scalasca.
